\section{Benchmark configurations and coverage}
\label{app:methods}

This appendix lists the configuration each method entered the benchmark with,
the runs excluded from it, and the reproducibility controls.
\Cref{sec:method} describes the tree-growing loop, and
\cref{app:complexity-details} gives its runtime and memory accounting.

\subsection{Selected configurations}
\label{app:selected-configs}

\Cref{tab:selected-configs} lists the configuration each method family entered
the benchmark with, after the selection rule of \cref{sec:experiments-repro}
picked one configuration per family by its mean real-data score. Candidates is
the size of the family's configuration grid. Settings not listed are the
library defaults. CIF uses stratified bootstrap sampling in classification and
ordinary bootstrap sampling in regression.

\begin{table}[!htbp]
  \centering
  \footnotesize
  \setlength{\tabcolsep}{4pt}
  \begin{tabular}{@{}llcp{7.4cm}@{}}
    \toprule
    \tablehead{Family} & \tablehead{Task} & Candidates & \tablehead{Selected configuration} \\
    \midrule
    CIT & Classification & 4 & selector MC; minimum budgets with adaptive stopping; histogram with 256 bins (257 candidates); Bonferroni; no honesty \\
    CIT & Regression & 4 & selector RDC; otherwise as classification \\
    CIF & Classification & 4 & selector RDC; 100 trees; stratified bootstrap; square-root feature sampling; no honesty; otherwise as CIT \\
    CIF & Regression & 4 & selector RDC; 100 trees; bootstrap; square-root feature sampling; honesty with a 0.5 split \\
    RF, ET & Both & 1 & 100 trees, library defaults, impurity importance \\
    XGBoost & Both & 5 & 500 rounds, depth 6, learning rate 0.1; importance total gain (classification) or total cover (regression) \\
    LightGBM & Both & 2 & 500 rounds, depth 6, learning rate 0.1; gain importance \\
    CatBoost & Both & 1 & 500 iterations, depth 6, learning rate 0.1 \\
    partykit ctree & Both & 2 & quadratic statistic, Monte Carlo test with 9,999 resamples \\
    partykit cforest & Both & 4 & as ctree; 100 trees, square-root mtry, unconditional out-of-bag permutation importance, partykit default (one permutation per tree over 100 trees) \\
    Permutation filters & Both & 1 & one permutation test per feature (MC, RDC; PC, DC, RDC) \\
    Boruta & Both & 1 & 200 iterations on an auto-sized random forest \\
    PI, CPI & Both & 1 & 10 permutation repeats on a random forest \\
    RF-RFE & Both & 1 & step 1 on a random forest \\
    \bottomrule
  \end{tabular}
  \caption{Configurations entered in the benchmark, one per family and task,
  with the size of the grid the selection rule chose from. The
  minimum budget is $\lceil 1/\alpha_{\mathrm{adj}} \rceil$ permutations
  at the Bonferroni-adjusted level; the histogram search with 256 bins offers at
  most 257 candidate thresholds per feature, so the per-threshold rule needs a
  budget of at least $\lceil 257/0.05 \rceil = 5{,}140$ to reject.}
  \label{tab:selected-configs}
\end{table}

\subsection{Benchmark coverage}

Each pairwise comparison uses the datasets, downstream learners, and standard
values of $k$ where both methods are present. Analyses that compare all
methods restrict to datasets with every method, learner, and standard $k$
present, which gives 13 classification and 6 regression complete-case
datasets. On \nolinkurl{gisette}, \nolinkurl{isolet}, and \nolinkurl{orlraws10P},
34 classification runs were excluded (\cref{tab:coverage-skips}). Fifteen
partykit ctree and cforest runs did not complete. Eighteen RDC-selector CIT and
CIF runs were lost to a host fault while the other seeds of the same cells
completed. One run on \nolinkurl{orlraws10P} was
excluded by author decision. No regression runs are excluded.

\begin{table}[!htbp]
  \centering
  \footnotesize
  \setlength{\tabcolsep}{4pt}
  \begin{tabular}{@{}llccl@{}}
    \toprule
    \tablehead{Dataset} & \tablehead{Method} & {Configurations} & {Seed runs} & \tablehead{Basis} \\
    \midrule
    \nolinkurl{gisette} & CIT with the RDC selector & 2 & 8 & host fault \\
    \nolinkurl{gisette} & R ctree & 1 & 5 & incomplete \\
    \nolinkurl{isolet} & CIF with the RDC selector & 2 & 4 & host fault \\
    \nolinkurl{isolet} & CIT with the RDC selector & 2 & 6 & host fault \\
    \nolinkurl{isolet} & R cforest & 1 & 5 & incomplete \\
    \nolinkurl{isolet} & R ctree & 1 & 5 & incomplete \\
    \nolinkurl{orlraws10P} & CIT with the RDC selector & 1 & 1 & author directed \\
    \bottomrule
  \end{tabular}
  \caption{Classification benchmark exclusions. Each seed run is one method
  configuration on one dataset with one seed; 34 runs are excluded in total,
  and all are operational exclusions. Incomplete rows are partykit ctree and
  cforest runs that did not complete; the 18 host-fault runs are RDC-selector
  CIT and CIF runs lost to a host fault; the
  single \nolinkurl{orlraws10P} run was excluded by author decision. No regression runs
  are excluded.}
  \label{tab:coverage-skips}
\end{table}

\subsection{Randomness and reproducibility}

The benchmark configuration fixes the seed schedule used for all reported
runs. Pure Python permutation routines use local generators. Numba parallel
permutation paths assign one seed per resample, which avoids sharing a mutable
generator across parallel loop iterations. The corrected R bridge was checked
by refitting one published ctree example inside R and through the harness with
identical splits; the seed schedule, rank directions, and the CPI threshold are
pinned by tests in the repository.
